\documentclass[11pt]{article}
%DIF LATEXDIFF DIFFERENCE FILE
%DIF DEL proposal-before.tex   Tue Mar 10 05:20:02 2026
%DIF ADD proposal.tex          Tue Mar 10 06:11:51 2026

% --- Packages ---
\usepackage[margin=1in]{geometry}
\usepackage{amsmath,amssymb}
\usepackage{graphicx}
\usepackage[hidelinks]{hyperref}
\usepackage{natbib}
\usepackage{booktabs}
\usepackage{enumitem}
\usepackage{xcolor}
\usepackage{siunitx}
\usepackage{subcaption}
\usepackage{titlesec}
\usepackage{setspace}

% --- Formatting ---
\titleformat{\section}{\normalfont\Large\bfseries}{\thesection}{1em}{}
\titleformat{\subsection}{\normalfont\large\bfseries}{\thesubsection}{1em}{}
\setlength{\parskip}{0.5em}
\setlength{\parindent}{0em}

% --- Draft helpers (remove for final) ---
\newcommand{\todo}[1]{\textcolor{red}{\textbf{[TODO: #1]}}}
\newcommand{\note}[1]{\textcolor{blue}{\textbf{[NOTE: #1]}}}

% =============================================================================
%DIF PREAMBLE EXTENSION ADDED BY LATEXDIFF
%DIF UNDERLINE PREAMBLE %DIF PREAMBLE
\RequirePackage[normalem]{ulem} %DIF PREAMBLE
\RequirePackage{color}\definecolor{RED}{rgb}{1,0,0}\definecolor{BLUE}{rgb}{0,0,1} %DIF PREAMBLE
\providecommand{\DIFaddtex}[1]{{\protect\color{blue}\uwave{#1}}} %DIF PREAMBLE
\providecommand{\DIFdeltex}[1]{{\protect\color{red}\sout{#1}}}                      %DIF PREAMBLE
%DIF SAFE PREAMBLE %DIF PREAMBLE
\providecommand{\DIFaddbegin}{} %DIF PREAMBLE
\providecommand{\DIFaddend}{} %DIF PREAMBLE
\providecommand{\DIFdelbegin}{} %DIF PREAMBLE
\providecommand{\DIFdelend}{} %DIF PREAMBLE
\providecommand{\DIFmodbegin}{} %DIF PREAMBLE
\providecommand{\DIFmodend}{} %DIF PREAMBLE
%DIF FLOATSAFE PREAMBLE %DIF PREAMBLE
\providecommand{\DIFaddFL}[1]{\DIFadd{#1}} %DIF PREAMBLE
\providecommand{\DIFdelFL}[1]{\DIFdel{#1}} %DIF PREAMBLE
\providecommand{\DIFaddbeginFL}{} %DIF PREAMBLE
\providecommand{\DIFaddendFL}{} %DIF PREAMBLE
\providecommand{\DIFdelbeginFL}{} %DIF PREAMBLE
\providecommand{\DIFdelendFL}{} %DIF PREAMBLE
%DIF HYPERREF PREAMBLE %DIF PREAMBLE
\providecommand{\DIFadd}[1]{\texorpdfstring{\DIFaddtex{#1}}{#1}} %DIF PREAMBLE
\providecommand{\DIFdel}[1]{\texorpdfstring{\DIFdeltex{#1}}{}} %DIF PREAMBLE
\newcommand{\DIFscaledelfig}{0.5}
%DIF HIGHLIGHTGRAPHICS PREAMBLE %DIF PREAMBLE
\RequirePackage{settobox} %DIF PREAMBLE
\RequirePackage{letltxmacro} %DIF PREAMBLE
\newsavebox{\DIFdelgraphicsbox} %DIF PREAMBLE
\newlength{\DIFdelgraphicswidth} %DIF PREAMBLE
\newlength{\DIFdelgraphicsheight} %DIF PREAMBLE
% store original definition of \includegraphics %DIF PREAMBLE
\LetLtxMacro{\DIFOincludegraphics}{\includegraphics} %DIF PREAMBLE
\newcommand{\DIFaddincludegraphics}[2][]{{\color{blue}\fbox{\DIFOincludegraphics[#1]{#2}}}} %DIF PREAMBLE
\newcommand{\DIFdelincludegraphics}[2][]{% %DIF PREAMBLE
\sbox{\DIFdelgraphicsbox}{\DIFOincludegraphics[#1]{#2}}% %DIF PREAMBLE
\settoboxwidth{\DIFdelgraphicswidth}{\DIFdelgraphicsbox} %DIF PREAMBLE
\settoboxtotalheight{\DIFdelgraphicsheight}{\DIFdelgraphicsbox} %DIF PREAMBLE
\scalebox{\DIFscaledelfig}{% %DIF PREAMBLE
\parbox[b]{\DIFdelgraphicswidth}{\usebox{\DIFdelgraphicsbox}\\[-\baselineskip] \rule{\DIFdelgraphicswidth}{0em}}\llap{\resizebox{\DIFdelgraphicswidth}{\DIFdelgraphicsheight}{% %DIF PREAMBLE
\setlength{\unitlength}{\DIFdelgraphicswidth}% %DIF PREAMBLE
\begin{picture}(1,1)% %DIF PREAMBLE
\thicklines\linethickness{2pt} %DIF PREAMBLE
{\color[rgb]{1,0,0}\put(0,0){\framebox(1,1){}}}% %DIF PREAMBLE
{\color[rgb]{1,0,0}\put(0,0){\line( 1,1){1}}}% %DIF PREAMBLE
{\color[rgb]{1,0,0}\put(0,1){\line(1,-1){1}}}% %DIF PREAMBLE
\end{picture}% %DIF PREAMBLE
}\hspace*{3pt}}} %DIF PREAMBLE
} %DIF PREAMBLE
\LetLtxMacro{\DIFOaddbegin}{\DIFaddbegin} %DIF PREAMBLE
\LetLtxMacro{\DIFOaddend}{\DIFaddend} %DIF PREAMBLE
\LetLtxMacro{\DIFOdelbegin}{\DIFdelbegin} %DIF PREAMBLE
\LetLtxMacro{\DIFOdelend}{\DIFdelend} %DIF PREAMBLE
\DeclareRobustCommand{\DIFaddbegin}{\DIFOaddbegin \let\includegraphics\DIFaddincludegraphics} %DIF PREAMBLE
\DeclareRobustCommand{\DIFaddend}{\DIFOaddend \let\includegraphics\DIFOincludegraphics} %DIF PREAMBLE
\DeclareRobustCommand{\DIFdelbegin}{\DIFOdelbegin \let\includegraphics\DIFdelincludegraphics} %DIF PREAMBLE
\DeclareRobustCommand{\DIFdelend}{\DIFOaddend \let\includegraphics\DIFOincludegraphics} %DIF PREAMBLE
\LetLtxMacro{\DIFOaddbeginFL}{\DIFaddbeginFL} %DIF PREAMBLE
\LetLtxMacro{\DIFOaddendFL}{\DIFaddendFL} %DIF PREAMBLE
\LetLtxMacro{\DIFOdelbeginFL}{\DIFdelbeginFL} %DIF PREAMBLE
\LetLtxMacro{\DIFOdelendFL}{\DIFdelendFL} %DIF PREAMBLE
\DeclareRobustCommand{\DIFaddbeginFL}{\DIFOaddbeginFL \let\includegraphics\DIFaddincludegraphics} %DIF PREAMBLE
\DeclareRobustCommand{\DIFaddendFL}{\DIFOaddendFL \let\includegraphics\DIFOincludegraphics} %DIF PREAMBLE
\DeclareRobustCommand{\DIFdelbeginFL}{\DIFOdelbeginFL \let\includegraphics\DIFdelincludegraphics} %DIF PREAMBLE
\DeclareRobustCommand{\DIFdelendFL}{\DIFOaddendFL \let\includegraphics\DIFOincludegraphics} %DIF PREAMBLE
%DIF COLORLISTINGS PREAMBLE %DIF PREAMBLE
\RequirePackage{listings} %DIF PREAMBLE
\RequirePackage{color} %DIF PREAMBLE
\lstdefinelanguage{DIFcode}{ %DIF PREAMBLE
%DIF DIFCODE_UNDERLINE %DIF PREAMBLE
  moredelim=[il][\color{red}\sout]{\%DIF\ <\ }, %DIF PREAMBLE
  moredelim=[il][\color{blue}\uwave]{\%DIF\ >\ } %DIF PREAMBLE
} %DIF PREAMBLE
\lstdefinestyle{DIFverbatimstyle}{ %DIF PREAMBLE
	language=DIFcode, %DIF PREAMBLE
	basicstyle=\ttfamily, %DIF PREAMBLE
	columns=fullflexible, %DIF PREAMBLE
	keepspaces=true %DIF PREAMBLE
} %DIF PREAMBLE
\lstnewenvironment{DIFverbatim}{\lstset{style=DIFverbatimstyle}}{} %DIF PREAMBLE
\lstnewenvironment{DIFverbatim*}{\lstset{style=DIFverbatimstyle,showspaces=true}}{} %DIF PREAMBLE
%DIF END PREAMBLE EXTENSION ADDED BY LATEXDIFF

\begin{document}

% --- Cover Page (does not count toward 5-page limit) ---
% sections/coverpage.tex
% Cover page for BYU Mentored Research Grant Proposal

\thispagestyle{empty}
\begingroup
\setlength{\parskip}{0.2em}
\titlespacing*{\subsection}{0pt}{0.6em}{0.2em}

\begin{center}
{\Large \textbf{Mentored Research Grant (MRG)}}
\end{center}

% --- Project Title ---
\subsection*{Project Title}
Bayesian Optimization of Multi-Material 3D-Printed Tensegrity Structures for Energy Absorption

% --- Principal Investigator(s) ---
\subsection*{Principal Investigator(s)}
\textbf{PI Name:} Jeffrey R.\ Hill \\
\textbf{PI Department:} Mechanical Engineering \\
\textbf{Co-PI:} Sterling G.\ Baird, Department of Mechanical Engineering

% --- Abstract ---
\subsection*{Abstract}

We propose a two-year mentored research program in which 2 undergraduate
students design and optimize multi-material 3D-printed tensegrity-inspired
structures for maximizing energy absorption while limiting permanent deformation.  Tensegrity-inspired
structures---assemblies of rigid compression members (PLA) and flexible tension
members (TPU)---offer tunable mechanical
responses and high strength-to-weight ratios.  Using a multifidelity Bayesian
optimization framework, students will combine finite element simulations
(low-fidelity) with physical experiments (high-fidelity) to efficiently navigate
the high-dimensional design space.  Students will gain hands-on experience in
CAD design, multi-material 3D printing, impact testing, computational modeling,
and data-driven optimization.  Expected outcomes include undergraduate
conference presentations (UCUR, ASME IDETC), a peer-reviewed journal
submission, and development of marketable technical skills.  This project
provides a rich mentored learning environment at the intersection of structural
mechanics, advanced manufacturing, and machine learning. Undergraduates will lead both the computational and experimental components of the project, progressing from structured training in Year 1 to independent research leadership and peer mentoring in Year 2.

% --- Number of Students ---
\subsection*{Number of Students to be Mentored}
\textbf{Number of Undergraduate Students:} 2 \\
\textbf{Number of Graduate Students:} 1 (co-mentor, in support of creating the mentoring environment)

% --- Budget ---
\subsection*{Budget}

\vspace{-0.5em}
\begin{center}
\begin{tabular}{@{}lrrr@{}}
\toprule
\textbf{Item} & \textbf{Year 1} & \textbf{Year 2} & \textbf{Total} \\
\midrule
Undergraduate wages    & \$9,000  & \$9,000  & \$18,000 \\
Graduate wages         & \$1,250  & \$1,250  & \$2,500  \\
Supplies               & \$2,250  & \$750    & \$3,000  \\
Travel                 & \$0      & \$1,500  & \$1,500  \\
\midrule
\textbf{Total}         & \textbf{\$12,500} & \textbf{\$12,500} & \textbf{\$25,000} \\
\bottomrule
\end{tabular}
\end{center}
\vspace{-0.5em}

% --- Relationship to External Funding ---
\subsection*{Relationship to External Funding}

The PIs do not currently hold external funding for tensegrity or architected
materials research.  This MRG provides the primary support for establishing a
mentored undergraduate research program in this area.  If successful, the
preliminary data may inform a future NSF proposal; however, the MRG is
self-contained and its outcomes do not depend on external funding.

\endgroup
\newpage


% =============================================================================
\section{Research Motivation and Overview}
% =============================================================================

We propose a two-year mentored research program in which 2--3 undergraduate
students design and optimize \textbf{multi-material 3D-printed
\DIFdelbegin \DIFdel{tensegrity
}\DIFdelend \DIFaddbegin \DIFadd{tensegrity-inspired }\DIFaddend structures} for energy absorption.  Tensegrity (tensional
integrity) structures---assemblies of rigid compression members suspended within
a continuous network of tension members---exhibit remarkable strength-to-weight
ratios and tunable mechanical responses, making them attractive for protective
equipment, packaging, and aerospace applications
\citep{skelton2009tensegrity, fraternali2015tensegrity}.
\DIFaddbegin \DIFadd{Recent work has shown that }\emph{\DIFadd{tensegrity-inspired}} \DIFadd{architectures can be
directly 3D-printed and retain desirable load-limiting behavior under impact
\mbox{%DIFAUXCMD
\citep{pajunen2019design}}\hskip0pt%DIFAUXCMD
.
}\DIFaddend 

Recent advances in multi-material 3D printing enable the fabrication of
\DIFdelbegin \DIFdel{tensegrity-like }\DIFdelend \DIFaddbegin \DIFadd{tensegrity-inspired }\DIFaddend structures on a single print platform, using PLA for rigid
struts and TPU for flexible tension elements
\DIFaddbegin \DIFadd{\mbox{%DIFAUXCMD
\citep{ye2023multimaterial, khatri2024energy}}\hskip0pt%DIFAUXCMD
}\DIFaddend .  This eliminates tedious manual
assembly and opens a rich, high-dimensional design space.  Navigating this space
efficiently requires intelligent experimental design; \textbf{multifidelity
Bayesian optimization} (BO) provides exactly this capability by combining
inexpensive simulations (low-fidelity) with physical experiments (high-fidelity)
to guide the search toward optimal configurations
\citep{mo2023accelerated, perdikaris2017nonlinear}.

\textbf{The primary objective of this project is to provide a rich,
interdisciplinary mentored research experience for undergraduates.}  Students
will gain hands-on experience in CAD design, multi-material 3D printing, impact
testing, computational modeling, and data-driven optimization---skills directly
transferable to graduate school and engineering careers.  Each student will
lead a clearly scoped research project, present at conferences, and contribute
to a peer-reviewed publication.

% --- Placeholder for overview figure ---
% TODO: Replace with actual figure before final submission
\begin{figure}[ht]
\centering
\fbox{\parbox{0.85\textwidth}{\centering\vspace{2em}
\textbf{Figure placeholder:} Overview of project approach---tensegrity unit
cell design (left), multi-material 3D printing of PLA struts and TPU tension
elements (center), and multifidelity Bayesian optimization loop connecting
simulations to experiments (right).
\vspace{2em}}}
\caption{Conceptual overview of the proposed research and mentoring framework.}
\label{fig:overview}
\end{figure}

% =============================================================================
\section{Background}
% =============================================================================

\textbf{Tensegrity structures} consist of isolated compression members (bars or
struts) held in stable equilibrium by a continuous network of tension members
(cables or strings) \citep{skelton2009tensegrity, sultan2009tensegrity}.  Their
defining characteristic---that no two rigid bars touch---yields lightweight
assemblies with surprising load-bearing capacity, resilience, and tunable
nonlinear responses \citep{amendola2014experimental, zhang2015tensegrity}.
\DIFdelbegin \DIFdel{Prior work has demonstrated }\DIFdelend \DIFaddbegin \DIFadd{Pajunen et al.\ \mbox{%DIFAUXCMD
\citep{pajunen2019design} }\hskip0pt%DIFAUXCMD
demonstrated that
}\emph{\DIFadd{tensegrity-inspired}} \DIFadd{unit cells can be directly 3D-printed and exhibit
load-limiting plateaus under drop-weight impact, confirming }\DIFaddend their potential for
\DIFdelbegin \DIFdel{impact protection
\mbox{%DIFAUXCMD
\citep{fraternali2015tensegrity}}\hskip0pt%DIFAUXCMD
, but design optimization remains largely
unexplored due to the high-dimensional parameter space (bar/cable lengths,
thicknesses, stiffnesses, topology, and geometry)}\DIFdelend \DIFaddbegin \DIFadd{energy absorption without requiring manual cable assembly.
Prior work on multi-material rigid--soft printing
\mbox{%DIFAUXCMD
\citep{ye2023multimaterial, khatri2024energy} }\hskip0pt%DIFAUXCMD
has shown that PLA--TPU
combinations printed on a single FDM platform can achieve tunable energy
absorption, though interface durability under repeated impact remains an
active area of investigation}\DIFaddend .

\textbf{Bayesian optimization} is a sequential, model-based strategy for
optimizing expensive-to-evaluate black-box functions
\citep{shahriari2016taking, frazier2018tutorial}.  A Gaussian process surrogate
model is fitted to observed data, and an acquisition function balances
exploration with exploitation.  BO has been applied to materials discovery
\citep{lookman2019active} and metamaterial design \citep{wang2022bayesian}.
Critically, Mo et al.\ \citep{mo2023accelerated} recently demonstrated
\textbf{multifidelity BO} for architected materials, using inexpensive
simulations as a low-fidelity source and physical experiments as high-fidelity
data.  We adopt this framework directly: finite element simulations serve as the
low-fidelity source, and 3D-printed tensegrity impact tests provide
high-fidelity validation.  \DIFaddbegin \DIFadd{An early calibration campaign will verify
low-fidelity--high-fidelity correlation before full optimization begins, since
MFBO effectiveness depends on adequate correlation between fidelity levels.
}\DIFaddend 

% =============================================================================
\section{Student Research Project~1: Simulation and Bayesian Optimization}
% =============================================================================

\textbf{Student profile:} One junior or senior undergraduate with coursework in
mechanics of materials or computational methods.

\subsection*{Scope and Deliverables}

This student will develop and run the computational pipeline:

\begin{enumerate}[nosep]
    \item \textbf{Parameterize a tensegrity unit cell} suitable for simulation
        and multi-material 3D printing, including strut geometry (PLA), tension
        element geometry (TPU), topology, and boundary conditions.
    \item \textbf{Build finite element models} using ABAQUS or open-source FEA
        tools (e.g., FEniCS) to predict force--displacement response and energy
        absorption under compaction and drop-test loading.  These simulations
        serve as the \emph{low-fidelity} source in the multifidelity framework.
        \DIFaddbegin \DIFadd{An initial calibration set will validate simulation--experiment
        correlation before full optimization begins.
    }\DIFaddend \item \textbf{Implement a multifidelity Bayesian optimization loop}
        (following Mo et al.\ \citep{mo2023accelerated}) using Python libraries
        (BoTorch/Ax) to efficiently explore the design space, fusing simulation
        predictions with experimental results from Student~2.
    \item \textbf{Identify optimal designs} and generate predictions for
        experimental validation.
\end{enumerate}

\textbf{Mentoring outcomes:} The student will develop skills in finite element
analysis, scientific computing (Python), and data-driven optimization.  They
will present results at \textbf{ASME IDETC} or \textbf{UCUR} and contribute as
co-author on a journal manuscript.

% =============================================================================
\section{Student Research Project~2: Fabrication, CAD, and Experimental Testing}
% =============================================================================

\textbf{Student profile:} One junior or senior undergraduate with coursework in
machine design, CAD, or compliant mechanisms.

\subsection*{Scope and Deliverables}

This student will lead fabrication and physical testing:

\begin{enumerate}[nosep]
    \item \textbf{Design \DIFdelbegin \DIFdel{tensegrity }\DIFdelend \DIFaddbegin \DIFadd{tensegrity-inspired }\DIFaddend structures in CAD} (SolidWorks or
        Fusion 360) for multi-material 3D printing, with PLA struts and TPU
        tension elements printed on a single platform\DIFaddbegin \DIFadd{, incorporating
        interface design strategies (e.g., interlocking geometries) to ensure
        durable PLA--TPU bonds \mbox{%DIFAUXCMD
\citep{ye2023multimaterial}}\hskip0pt%DIFAUXCMD
}\DIFaddend .
    \item \textbf{Fabricate structures} using the department's multi-material
        3D printer, iterating on print parameters to achieve reliable
        tensegrity geometries.
    \item \textbf{Conduct experimental testing:} compaction tests (quasi-static
        compression), drop tests (dynamic impact), and wave propagation
        measurements (hitting on one end and measuring time response on the
        other).  These experiments serve as the \emph{high-fidelity} source.
    \item \textbf{Compare experimental results} against simulation predictions
        from Student~1, quantifying model accuracy and informing surrogate
        updates.
\end{enumerate}

\textbf{Mentoring outcomes:} The student will develop skills in CAD, additive
manufacturing, test design, and data analysis.  They will present at
\textbf{UCUR} or \textbf{ASME IDETC} and contribute as co-author on a journal
manuscript.

% =============================================================================
\section{Mentoring Environment}
\label{sec:mentoring}
% =============================================================================

\textbf{The mentoring of undergraduate researchers is the central objective of
this proposal.}  We are committed to providing a structured, supportive
environment that develops students into independent, capable researchers.

\subsection*{Recruitment}

Students will be recruited from courses in mechanics of materials, machine
design, kinematics, and compliant mechanisms through brief in-class
presentations.  We target junior or senior undergraduates.  We will actively
seek to recruit students from underrepresented groups in engineering.

\subsection*{Mentoring Structure}

\begin{itemize}[nosep]
    \item \textbf{Weekly individual meetings} (30 min) with faculty PI
        (Prof.~Hill) and Co-PI (Prof.~Baird) for technical guidance,
        troubleshooting, and professional development discussions.
    \item \textbf{Weekly lab group meetings} (1 hour) including: progress
        updates, student presentations on recent literature, career
        discussions, and skill-building workshops (e.g., scientific writing,
        conference presentation skills, research ethics).
    \item \textbf{Graduate student co-mentor:} Prof.~Baird's master's student
        will serve as day-to-day co-mentor, providing hands-on guidance in
        simulation, coding, fabrication, and experimental technique.  This
        layered mentoring model ensures students receive timely support while
        the graduate student develops their own mentoring skills.
    \item \textbf{Progressive responsibility:} Students begin with structured
        tasks (running simulations with provided scripts, assembling structures
        from detailed CAD plans) and progress to independent design decisions,
        culminating in ownership of their research project.
    \item \textbf{Peer mentoring:} In Year~2, experienced students will mentor
        newly recruited freshmen or sophomores, reinforcing their own learning
        while expanding the research group and creating a sustainable pipeline
        of trained researchers.
\end{itemize}

\subsection*{Student Development Outcomes}

Students will develop a broad, transferable skill set spanning:
\begin{itemize}[nosep]
    \item \textbf{Technical:} experimental mechanics (test design, data
        acquisition), computation (FEA modeling, Python, Bayesian
        optimization), and design (CAD, multi-material 3D printing).
    \item \textbf{Communication:} lab notebooks, group presentations,
        conference talks, and co-authoring a peer-reviewed manuscript.
    \item \textbf{Professional:} project management, collaboration with
        interdisciplinary teams, and preparation for graduate school or
        industry careers.
\end{itemize}

% =============================================================================
\section{Expected Research Outcomes}
% =============================================================================

\begin{enumerate}[nosep]
    \item \textbf{Trained undergraduate researchers:} 2--3 students with
        hands-on experience in advanced manufacturing, computational modeling,
        and data-driven optimization, prepared for graduate study or engineering
        careers.
    \item \textbf{Conference presentations:} Student-led presentations at
        \textbf{UCUR} (Utah Conference on Undergraduate Research) and
        \textbf{ASME IDETC} (International Design Engineering Technical
        Conferences).
    \item \textbf{Journal submission:} One peer-reviewed manuscript submitted to
        the \emph{ASME Journal of Mechanical Design} or \emph{Smart Materials
        and Structures}, with undergraduate students as co-authors.
    \item \textbf{Open-source tools:} Simulation code, optimization scripts,
        and experimental data published on GitHub for reproducibility.
\end{enumerate}

% =============================================================================
\section{Potential Impact of Work}
% =============================================================================

This project advances the design of architected materials for energy absorption
with direct applications in \textbf{protective equipment} (helmets, body
armor), \textbf{packaging} (shock-absorbing inserts), and \textbf{aerospace}
(lightweight impact-resistant structures).  By demonstrating that multifidelity
Bayesian optimization can efficiently navigate the tensegrity design space, we
establish a generalizable framework applicable to other classes of architected
materials and metamaterials.

The interdisciplinary training---spanning structural mechanics, advanced
manufacturing, and machine learning---prepares students for graduate study or
careers in industries increasingly reliant on data-driven design.  The
open-source dissemination of tools and data extends impact beyond BYU and
enables other research groups to build on our work.

% =============================================================================
\section{Timeline}
% =============================================================================

\begin{table}[ht]
\centering
\small
\begin{tabular}{@{}lcccc@{}}
\toprule
\textbf{Task} & \textbf{Y1 Fall} & \textbf{Y1 Winter} & \textbf{Y2 Fall} & \textbf{Y2 Winter} \\
\midrule
Student recruitment \& onboarding    & \(\bullet\) &           &           &           \\
Literature review                     & \(\bullet\) & \(\bullet\) &           &           \\
Tensegrity parameterization \& CAD   & \(\bullet\) & \(\bullet\) &           &           \\
FEA simulation development           &             & \(\bullet\) & \(\bullet\) &           \\
Multifidelity BO implementation      &             & \(\bullet\) & \(\bullet\) &           \\
Multi-material 3D printing \& fabrication &        & \(\bullet\) & \(\bullet\) & \(\bullet\) \\
Impact testing \& data collection    &             &             & \(\bullet\) & \(\bullet\) \\
Simulation--experiment comparison    &             &             &             & \(\bullet\) \\
Conference presentations (UCUR/IDETC)&             &             & \(\bullet\) & \(\bullet\) \\
Journal manuscript preparation       &             &             &             & \(\bullet\) \\
Peer mentoring of new students       &             &             & \(\bullet\) & \(\bullet\) \\
\bottomrule
\end{tabular}
\caption{Project timeline across four semesters.}
\label{tab:timeline}
\end{table}

% =============================================================================
\section{Budget}
% =============================================================================

% sections/budget.tex
% Budget justification for BYU Mentored Research Grant

Total requested: \textbf{\$25,000} over 2 years.  Per grant guidelines, no
funds are allocated to equipment, tuition, or faculty wages.

\begin{center}
{\small
\begin{tabular}{@{}llrr@{}}
\toprule
\textbf{Category} & \textbf{Description} & \textbf{Year 1} & \textbf{Year 2} \\
\midrule
Undergraduate wages
    & 2 students, part-time + summer full-time
    & \$9,000 & \$9,000 \\
Graduate student wages
    & Partial RA for grad co-mentor
    & \$1,250 & \$1,250 \\
Supplies
    & Filament, load cells, DAQ, compute
    & \$2,000 & \$500 \\
Travel
    & Conference registration \& travel
    & \$0 & \$1,500 \\
Other
    & Poster printing, consumables
    & \$250 & \$250 \\
\midrule
\textbf{Annual Total} & & \textbf{\$12,500} & \textbf{\$12,500} \\
\bottomrule
\end{tabular}
}
\end{center}

\textbf{Undergraduate wages (\$18,000):} Supports 2 undergraduates working
part-time (\textasciitilde 10~hrs/week) during fall and winter semesters and
full-time (\textasciitilde 40~hrs/week) during summer terms for simulation,
3D printing, and testing.
%
\textbf{Graduate student wages (\$2,500):} Prof.~Baird's master's student
serves as day-to-day co-mentor.
%
\textbf{Supplies (\$2,500):} PLA/TPU filament, load cells, DAQ module,
fasteners, and cloud compute credits (Year~1 heavy).
%
\textbf{Travel (\$1,500):} Student presentations at ASME IDETC or UCUR
(Year~2).
%
\textbf{Other (\$500):} Poster printing, consumables, contingency.


% =============================================================================
% --- References (do not count toward 5-page limit) ---
\bibliographystyle{plainnat}
\bibliography{references}

% --- Bio Sketches (do not count toward 5-page limit) ---
\newpage
% sections/biosketch.tex
% Biographical Sketches for PI and Co-PI (NIH Biosketch format)

\section*{Biographical Sketches}

% ============================================================================
\subsection*{Jeffrey R.\ Hill --- Principal Investigator}
% ============================================================================

\subsubsection*{A.\ Personal Statement}

I am an Associate Professor of Mechanical Engineering at Brigham Young
University with expertise in tensegrity structures, smart materials, and
experimental mechanics.  My research group has recently published on tensegrity
robot locomotion, cable tension interactions, and 3D-printed tension networks,
providing direct expertise relevant to this proposal.  I have a strong record
of mentoring undergraduate and graduate students through hands-on research,
which aligns with the goals of the Mentored Research Grant program.

\begin{enumerate}[nosep]
    \item Layer B, Denning H, Hill J. Control and locomotion of tensegrity robots through manipulation of the center of mass. \emph{Robotica}. 2024. \url{https://doi.org/10.1017/S0263574724001139}
    \item Masmeijer T, Swain C, Hill J, Habtour E. Programming tension in 3D printed networks inspired by spiderwebs. \emph{Materials \& Design}. 2025. \url{https://doi.org/10.1016/j.matdes.2025.115281}
\end{enumerate}

\subsubsection*{B.\ Positions, Scientific Appointments, and Honors}

\textbf{Education/Training}

\begin{tabular}{@{}llll@{}}
\textbf{Institution} & \textbf{Degree} & \textbf{Completion Date} & \textbf{Field of Study} \\
\hline
Pennsylvania State University & Ph.D. & 2011 & Mechanical Engineering \\
Brigham Young University & M.S. & 2005 & Mechanical Engineering \\
Brigham Young University & B.S. & 2004 & Mechanical Engineering \\
\end{tabular}

\textbf{Positions and Employment}

\begin{itemize}[nosep]
    \item 2021--present: Associate Professor, Dept.\ of Mechanical Engineering, Brigham Young University
    \item 2016--2021: Principal Member of the Technical Staff, Sandia National Laboratories
    \item 2011--2016: Senior Member of the Technical Staff, Sandia National Laboratories
    \item 2008--2011: Visiting Research Investigator, University of Michigan
\end{itemize}

\subsubsection*{C.\ Contributions to Science}

\textbf{1.\ Tensegrity structures and robots.}
My group investigates design, construction, and control of tensegrity systems.
We have demonstrated simplified construction methods, characterized cable
tension interactions, and developed locomotion strategies for tensegrity
robots.

\begin{enumerate}[nosep]
    \item Layer B, Denning H, Hill J. Control and locomotion of tensegrity robots through manipulation of the center of mass. \emph{Robotica}. 2024. \url{https://doi.org/10.1017/S0263574724001139}
    \item Denning H, Bullinger A, Hill J. Tuning Tensegrity: Simplified Construction and Cable Tension Interactions. \emph{ASME SMASIS 2024}. \url{https://doi.org/10.1115/SMASIS2024-140381}
    \item Brown A, Swain C, Usevitch N, Hill J. Development of a Continuous Cable Tensegrity. \emph{ASME SMASIS 2024}. \url{https://doi.org/10.1115/SMASIS2024-140185}
\end{enumerate}

\textbf{2.\ Advanced manufacturing and structural design.}
I have contributed to 3D-printed tension networks inspired by spiderwebs and
encapsulation techniques for protecting electronic components, demonstrating
expertise in additive manufacturing and structural characterization.

\begin{enumerate}[nosep]
    \item Masmeijer T, Swain C, Hill J, Habtour E. Programming tension in 3D printed networks inspired by spiderwebs. \emph{Materials \& Design}. 2025. \url{https://doi.org/10.1016/j.matdes.2025.115281}
    \item Hill J, Chen A, Wilson N, Boll C, O'Bannon M, Trentman D. Microbead Encapsulation for Protection of Electronic Components. \emph{J.\ Electronic Packaging}. 2025;147(021007). \url{https://doi.org/10.1115/1.4067650}
\end{enumerate}

\textbf{3.\ Smart materials and structural optimization.}
My earlier work focused on actuator grouping optimization and piezoelectric
fiber composite analysis for adaptive structures, establishing a foundation in
computational and experimental approaches to structural systems.

\begin{enumerate}[nosep]
    \item Hill J, Wang K. Development of the En Masse Elimination Algorithm for Actuator Grouping Optimization. \emph{J.\ Intelligent Material Systems and Structures}. 2011;22(11):1203--1212. \url{https://doi.org/10.1177/1045389X11417651}
    \item Kim J, Hill J, Wang K. An asymptotic approach for the analysis of piezoelectric fiber composite beams. \emph{Smart Materials and Structures}. 2011;20(2):025023.
    \item Hill J, Wang K, Fang H, Quijano U. Actuator grouping optimization on flexible space reflectors. \emph{SPIE Active and Passive Smart Structures}. 2011. \url{https://doi.org/10.1117/12.880086}
\end{enumerate}

\vspace{2em}

% ============================================================================
\subsection*{Sterling G.\ Baird --- Co-Principal Investigator}
% ============================================================================

\subsubsection*{A.\ Personal Statement}

I am an Assistant Professor of Mechanical Engineering at Brigham Young
University specializing in Bayesian optimization, self-driving laboratories,
and data-driven materials discovery.  My work on the Honegumi interface and
high-dimensional Bayesian optimization directly supports the optimization
component of this proposal.  I have extensive experience mentoring
undergraduate researchers and developing accessible open-source tools that
lower barriers to adopting advanced optimization methods.

\begin{enumerate}[nosep]
    \item Baird SG, Falkowski AR, Sparks TD. Honegumi: An Interface for Accelerating the Adoption of Bayesian Optimization in the Experimental Sciences. 2025. \url{https://doi.org/10.48550/arXiv.2502.06815}
    \item Tom G, Schmid SP, Baird SG, et al. Self-Driving Laboratories for Chemistry and Materials Science. \emph{Chemical Reviews}. 2024;124(16):9633--9732. \url{https://doi.org/10.1021/acs.chemrev.4c00055}
\end{enumerate}

\subsubsection*{B.\ Positions, Scientific Appointments, and Honors}

\textbf{Education/Training}

\begin{tabular}{@{}llll@{}}
\textbf{Institution} & \textbf{Degree} & \textbf{Completion Date} & \textbf{Field of Study} \\
\hline
University of Utah & Ph.D. & 2023 & Materials Science and Engineering \\
Brigham Young University & M.S. & 2021 & Mechanical Engineering \\
Brigham Young University & B.S. & 2018 & Applied Physics \\
\end{tabular}

\textbf{Positions and Employment}

\begin{itemize}[nosep]
    \item 2025--present: Assistant Professor, Dept.\ of Mechanical Engineering, Brigham Young University
    \item 2025: Principal Investigator and Staff Scientist, Acceleration Consortium, University of Toronto
    \item 2023--2025: Director of Training and Programs, Acceleration Consortium, University of Toronto
\end{itemize}

\textbf{Honors}
\begin{itemize}[nosep]
    \item Guest Editor, \emph{npj Computational Materials}
    \item Journal reviewer for \emph{Nature Communications}, \emph{Digital Discovery}, \emph{JOSS}
\end{itemize}

\subsubsection*{C.\ Contributions to Science}

\textbf{1.\ Bayesian optimization for materials and experimental sciences.}
I developed Honegumi, an interactive interface for generating Bayesian
optimization scripts, and demonstrated high-dimensional BO applied to
materials property prediction and formulation optimization.

\begin{enumerate}[nosep]
    \item Baird SG, Falkowski AR, Sparks TD. Honegumi: An Interface for Accelerating the Adoption of Bayesian Optimization in the Experimental Sciences. 2025. \url{https://doi.org/10.48550/arXiv.2502.06815}
    \item Baird SG, Liu M, Sparks TD. High-Dimensional Bayesian Optimization of 23 Hyperparameters over 100 Iterations for an Attention-Based Network to Predict Materials Property. \emph{Computational Materials Science}. 2022;211:111505. \url{https://doi.org/10.1016/j.commatsci.2022.111505}
    \item Baird SG, Hall JR, Sparks TD. Compactness Matters: Improving Bayesian Optimization Efficiency of Materials Formulations through Invariant Search Spaces. \emph{Computational Materials Science}. 2023;224:112134. \url{https://doi.org/10.1016/j.commatsci.2023.112134}
    \item Baird SG, et al. Bayesian Optimization Hackathon for Chemistry and Materials. 2025. \url{https://doi.org/10.26434/chemrxiv-2025-dzh5z}
\end{enumerate}

\textbf{2.\ Self-driving laboratories.}
I co-authored the definitive review of self-driving labs for chemistry and
materials science and developed a low-cost ``Hello World'' demonstration
platform, making autonomous experimentation accessible to a broad audience.

\begin{enumerate}[nosep]
    \item Tom G, Schmid SP, Baird SG, et al. Self-Driving Laboratories for Chemistry and Materials Science. \emph{Chemical Reviews}. 2024;124(16):9633--9732. \url{https://doi.org/10.1021/acs.chemrev.4c00055}
    \item Baird SG, Sparks TD. Building a ``Hello World'' for Self-Driving Labs: The Closed-loop Spectroscopy Lab Light-mixing Demo. \emph{STAR Protocols}. 2023;4(2):102329. \url{https://doi.org/10.1016/j.xpro.2023.102329}
    \item Lo S, Baird SG, et al. Review of Low-Cost Self-Driving Laboratories in Chemistry and Materials Science: The ``Frugal Twin'' Concept. \emph{Digital Discovery}. 2024;3(5):842--868. \url{https://doi.org/10.1039/D3DD00223C}
\end{enumerate}

\textbf{3.\ Community building and open-source tools for data-driven research.}
I organized a Bayesian optimization hackathon (120 participants, 62
organizations, 14 countries) and developed open-source tools that lower the
barrier to adopting data-driven methods in experimental research.  I have
mentored 8 undergraduate research assistants at BYU and 12+ staff and students
at the Acceleration Consortium.

\begin{enumerate}[nosep]
    \item Baird SG, et al. Bayesian Optimization Hackathon for Chemistry and Materials. 2025. \url{https://doi.org/10.26434/chemrxiv-2025-dzh5z}
    \item Seegmiller CC, Baird SG, Sayeed HM, Sparks TD. Discovering Chemically Novel, High-Temperature Superconductors. \emph{Computational Materials Science}. 2023.
\end{enumerate}


\end{document}
